\documentclass[11pt,a4paper]{article}

\usepackage[T1]{fontenc}
\usepackage[utf8]{inputenc}
\usepackage{microtype}
\usepackage{booktabs}
\usepackage{multirow}
\usepackage{amsmath}
\usepackage{amssymb}
\usepackage{xcolor}
\usepackage{listings}
\usepackage{hyperref}
\usepackage{natbib}
\usepackage{geometry}
\geometry{margin=2.5cm}

\lstset{
  basicstyle=\ttfamily\small,
  breaklines=true,
  frame=single,
  columns=fullflexible,
}

\title{Linking a Conceptual Metaphor Thesaurus to WordNet\\
       via Definition Similarity and Hypernym Chain Matching}

\author{Francis Bond \and [co-authors]}

\date{\today}

\begin{document}
\maketitle

\begin{abstract}
We describe an automatic method for linking the entries of a large
English conceptual metaphor thesaurus to senses in the Open Multilingual
Wordnet \citep{bond-foster-2013-linking}.
Each thesaurus entry carries both a literal meaning and a metaphorical
meaning, together with the conceptual metaphor theme (e.g.\
\textsc{colour is mineral}) that situates the entry in the
\citealt{lakoff-johnson-1980} tradition.
We exploit this structure with two complementary scoring signals:
(i)~\emph{definition similarity}, computed either as Jaccard overlap on
content words or as cosine similarity of dense embeddings produced by a
locally hosted neural encoder (\texttt{embeddinggemma}); and
(ii)~\emph{hypernym chain matching}, which walks the WordNet hypernym
closure of each candidate synset and checks whether the \textsc{source}
or \textsc{target} domain keywords from the metaphor theme appear among
the ancestor lemmas.
For each entry we independently select the best-scoring sense for the
literal meaning and for the metaphorical meaning.
We present a running example (\emph{sapphire} in \textsc{colour is mineral}),
report preliminary results over the full thesaurus of 8,747~entries, and
discuss threshold selection and directions for refinement.
\end{abstract}

%% -------------------------------------------------------------------------
\section{Introduction}
\label{sec:intro}

Conceptual Metaphor Theory \citep{lakoff-johnson-1980} holds that
abstract concepts are systematically understood in terms of more concrete
domains through cross-domain mappings.
A classic example is \textsc{argument is war}, where reasoning about
arguments borrows vocabulary from the domain of warfare
(\emph{attack}, \emph{defend}, \emph{shoot down}).
Large manually curated resources that catalogue such mappings offer a
valuable window into the systematic nature of metaphorical language, but
they remain largely disconnected from the formal lexical-semantic
infrastructure of wordnets.

The thesaurus of \citet{hcp78} (hereafter \textsc{MetThes}) provides an
extensive lexical catalogue of English conceptual metaphors.
Each entry pairs a headword (a lexical item used metaphorically) with a
short gloss of its literal meaning, a statement of its metaphorical
meaning, and an illustrative sentence, all situated under a named
conceptual metaphor theme in the form \textsc{target is source}
(e.g.\ \textsc{quality is money/wealth}).
Connecting \textsc{MetThes} entries to WordNet senses would:
\begin{itemize}
  \item ground the literal meanings in a well-studied ontological hierarchy;
  \item make the source and target conceptual domains navigable via
        hypernym chains; and
  \item enable downstream lexical-semantic analyses (e.g.\ polysemy
        profiling, cross-lingual projection) that rely on synset identifiers.
\end{itemize}

In this paper we describe a fully automatic pipeline that, for each entry,
selects the most plausible WordNet sense for the literal meaning and a
(typically different) sense for the metaphorical meaning, combining
definition-level text similarity with ontological signal from the
hypernym graph.

%% -------------------------------------------------------------------------
\section{Resource Description}
\label{sec:resource}

\subsection{The Metaphor Thesaurus}

\textsc{MetThes} \citep{hcp78} is organised hierarchically into six
thematic parts (Table~\ref{tab:parts}), 339~conceptual metaphor themes,
and 8,747~entries across 6,880~distinct headwords.

\begin{table}[ht]
\centering
\caption{The six parts of \textsc{MetThes} and their entry counts.}
\label{tab:parts}
\begin{tabular}{llr}
\toprule
Part & Title & Entries \\
\midrule
I   & Values, Qualities and Quantities          & 1,430 \\
II  & Emotions, Experiences and Relationships   & 1,280 \\
III & Thinking \& Communicating                 & 1,896 \\
IV  & Activities, Processes, Movement and Time  & 2,381 \\
V   & Humans, Society, Animals and Senses       & 1,391 \\
VI  & Things \& Substances                      &   369 \\
\midrule
    & \textbf{Total}                            & \textbf{8,747} \\
\bottomrule
\end{tabular}
\end{table}

Each theme is named in the form \textsc{target is source} (or
\textsc{target is source\textsubscript{1}/source\textsubscript{2}} for
multiple sources), yielding 259~distinct target domains and 298~distinct
source domains.
Entries within a theme are grouped into subsections with underlined
headings that provide more specific characterisations of the mapping.

A single entry contains the following fields:

\begin{description}
  \item[headword] The lexical item used metaphorically, e.g.\ \emph{sapphire}.
  \item[literal\_meaning] A short gloss of the word's literal sense,
    e.g.\ \emph{transparent bright blue precious stone}.
  \item[word\_class] Part of speech in literal and metaphorical use,
    drawn from a 22-tag inventory (\texttt{n}, \texttt{adj}, \texttt{v},
    \texttt{idi}, \ldots).
  \item[metaphorical\_meaning] The abstract concept the headword
    denotes in the metaphorical usage, written in small capitals,
    e.g.\ \textsc{bright colour}.
  \item[example] An illustrative sentence, e.g.\ \emph{she has sapphire
    blue eyes}.
\end{description}

\subsection{WordNet and the Open Multilingual Wordnet}

We use Princeton WordNet~3.1 as distributed via the Open Multilingual
Wordnet~2.0 (OMW\,2.0; \citealt{bond-foster-2013-linking}), accessed
through the \texttt{wn}~1.0 Python library.
OMW\,2.0 provides synset identifiers of the form
\texttt{omw-en-\textit{offset}-\textit{pos}}, together with definitions,
usage examples, lemmas, and hypernym/hyponym relations.
The English portion covers nouns (\texttt{n}), verbs (\texttt{v}),
adjectives (\texttt{a} and satellite adjectives \texttt{s}, merged to \texttt{a}), and adverbs (\texttt{r}).

%% -------------------------------------------------------------------------
\section{Method}
\label{sec:method}

For each thesaurus entry we retrieve all WordNet synsets whose lemmas
include the headword, regardless of part of speech.
We do not filter by POS at retrieval time because the POS recorded in the
thesaurus reflects the word's \emph{metaphorical} usage, which may differ
from the POS of its dominant literal WordNet senses (e.g.\ \emph{sapphire}
is tagged \texttt{adj} in the thesaurus because of the attributive use
\emph{sapphire blue}, yet its relevant literal senses are nouns).
Each candidate synset is then scored independently against the literal
meaning and the metaphorical meaning using two complementary signals.

\subsection{Definition Similarity}
\label{subsec:defsim}

We combine the synset definition and any usage examples into a single
string $d_s$ and compare it to the query string $q$ (the literal or
metaphorical meaning).

\paragraph{Jaccard overlap (baseline).}
Both strings are lowercased, tokenised on word boundaries, and filtered
against a 22-word stoplist.
The similarity is the Jaccard coefficient of the resulting content-word sets:
\begin{equation}
  \mathrm{sim}_{\mathrm{Jac}}(q, d_s)
    = \frac{|\mathit{tok}(q) \cap \mathit{tok}(d_s)|}
           {|\mathit{tok}(q) \cup \mathit{tok}(d_s)|}.
\end{equation}
This is fast and fully interpretable but fails when synonymous words are
used (e.g.\ \textsc{bright colour} vs.\ \emph{a light shade of blue}).

\paragraph{Embedding cosine similarity.}
To capture semantic equivalence beyond surface overlap we encode both
strings and compute cosine similarity:
\begin{equation}
  \mathrm{sim}_{\mathrm{emb}}(q, d_s)
    = \frac{\mathbf{e}_q \cdot \mathbf{e}_s}
           {\|\mathbf{e}_q\|\,\|\mathbf{e}_s\|},
\end{equation}
where $\mathbf{e}_q$ and $\mathbf{e}_s$ are the respective embedding vectors.
We support two encoders: \texttt{embeddinggemma} (768-dimensional, served
locally via Ollama) and \texttt{sup-simcse-roberta-large}~\cite{gao2021simcse}
(1024-dimensional, loaded via \texttt{sentence-transformers}), referred to as
SimCSE.
SimCSE is trained with supervised contrastive learning on natural language
inference data, making its representations sensitive to fine-grained meaning
differences — an advantage for distinguishing near-synonymous WordNet senses.

\paragraph{Batch encoding.}
All unique strings to be encoded (synset texts, literal meanings, and
metaphorical meanings across the entire thesaurus) are collected in a
first pass and sent to Ollama in batches of 64.
Results are persisted to a local cache keyed by model name and text, so
subsequent runs only encode new strings.
For the full thesaurus this amounts to approximately 31,600~unique texts
spanning roughly 21,950~synsets across 6,880~headwords.

\subsection{Hypernym Chain Matching}
\label{subsec:hypernym}

The metaphor theme name encodes the conceptual mapping directly.
For a theme of the form \textsc{target is source} we define:
\begin{itemize}
  \item \textbf{source domains} — the words to the right of \textsc{is},
    e.g.\ \{MINERAL\} for \textsc{colour is mineral};
  \item \textbf{target domains} — the words to the left of \textsc{is},
    e.g.\ \{COLOUR\} for \textsc{colour is mineral}.
\end{itemize}
Multiple domains separated by ``/'' are treated as a list (all must appear
for a full score).

For each candidate synset $s$ we perform a breadth-first search (BFS)
over the full hypernym closure (up to 300~visited synsets) and record the
\emph{first} (most specific) ancestor synset whose lemmas contain a
domain keyword.
The hypernym score is the fraction of domain keywords matched:
\begin{equation}
  h(s, D) = \frac{|\{d \in D \mid \exists\,\text{ancestor of }s
                   \text{ with lemma } \supseteq d\}|}{|D|}.
\end{equation}
Crucially, we record the matching ancestor synset identifier and definition,
so the result carries a grounded annotation such as:
\begin{center}
  \texttt{MINERAL} $\to$ \texttt{omw-en-14662574-n}
  (\emph{solid homogeneous inorganic substances \ldots}).
\end{center}
The hypernym signal is applied asymmetrically:
\begin{itemize}
  \item the \textbf{literal sense} is matched against the
        \textbf{source domains}
        (the word denotes something \emph{in} the source domain literally);
  \item the \textbf{metaphorical sense} is matched against the
        \textbf{target domains}
        (the word is used to talk about something \emph{in} the target domain
        metaphorically).
\end{itemize}

\subsection{Source Domain Taxonomy}
\label{subsec:domainmap}

The hypernym chain matching produces, for each entry, a record of which
ancestor synset triggered the domain match (e.g.\ MINERAL $\to$
\texttt{omw-en-14662574-n}).
Aggregating these records across all entries yields a \emph{source domain
map}: a mapping from each source domain keyword to the canonical WordNet
synset(s) that best represent it.

We build this map as follows.
For each source domain $D$ we collect all synsets that were reached via
hypernym matching, count how many entries triggered each synset, and group
them by WordNet part of speech.
Within each POS group we retain synsets matched at least twice (the
\texttt{--min-count} threshold) and treat each qualifying synset as a
distinct canonical representative of the domain in that POS.
Domains with a single canonical synset are represented as a leaf node;
domains with multiple canonical synsets are represented as a grouping node
with labelled children (e.g.\ \textsc{bird-1} through \textsc{bird-5}
for the five bird-order synsets identified for the domain \textsc{bird}).

The resulting taxonomy is stored as a flat TOML file in which every entry
carries a \texttt{name}, \texttt{pos}, \texttt{synset\_id}, \texttt{hypernym}
(the parent domain name), and \texttt{definition}.
Four virtual root nodes (\textsc{n\_top}, \textsc{v\_top}, \textsc{a\_top},
\textsc{r\_top}) serve as POS-level anchors.

\paragraph{Intra-domain WordNet ancestry.}
We further enrich the taxonomy by checking whether any single-synset
domain's canonical synset is a WordNet hyponym of another domain's synset.
When such a relationship holds, the more specific domain is promoted under
the broader one (e.g.\ \textsc{waterbird} $\to$ \textsc{bird}).
Of the 213 source domains with hypernym matches, 92 had their parent
updated in this way, producing a coherent shallow ontology of source domains.

\paragraph{Target domain map.}
We apply the same procedure to the metaphorical role, collecting hypernym-match
evidence from \texttt{wn\_metaphorical} scores to build a parallel
\emph{target} domain map.
The target map yields 128 domains (186 synset-bearing entries, 21 grouping nodes,
74 ancestry promotions), substantially fewer than the source map because
target domains in the thesaurus are drawn from a smaller lexical range
(abstract concepts such as \textsc{love}, \textsc{time}, \textsc{society}).

\paragraph{Combined domain map.}
We also produce a \emph{combined} map by summing evidence across both roles
before applying the frequency threshold.
This gives 325 domains (595 synset-bearing entries, 75 grouping nodes,
218 ancestry promotions) and is the most comprehensive taxonomy of the
concepts covered by the thesaurus.

\paragraph{Source vs target synset divergence.}
Of the 213 source and 128 target domains, only 16 names appear in both maps.
Of these, 4 have identical canonical synsets in both roles
(\textsc{animal}, \textsc{building}, \textsc{human}, \textsc{money}),
while 12 diverge.
The divergences reveal genuine polysemy: the same domain word attracts
different senses depending on whether it anchors a literal or a metaphorical
meaning.
For example, \textsc{food} maps to \textit{``any substance that can be metabolized''}
as a source domain but to \textit{``a substance that can be used or prepared as food''}
as a target domain; \textsc{race} maps to \textit{``a contest of speed''} (source)
versus \textit{``people who are believed to belong to a race''} (target);
and \textsc{plant} resolves to a biological cell unit (source) versus a woody
plant (target).
This pattern — source roles attracting the more general/process sense and
target roles attracting the more prototypical/object sense — is consistent with
the structure of conceptual metaphor, where abstract targets are understood
through concrete, well-delineated source concepts.

\subsection{Combined Score and Sense Selection}
\label{subsec:combined}

The two signals are combined with a weighting parameter $\alpha \in [0,1]$:
\begin{equation}
  \mathrm{score}(s, q, D)
    = \alpha \cdot \mathrm{sim}(q, d_s)
    + (1 - \alpha) \cdot h(s, D),
\end{equation}
where $\mathrm{sim}$ is either $\mathrm{sim}_{\mathrm{Jac}}$ or
$\mathrm{sim}_{\mathrm{emb}}$.
We use $\alpha = 0.5$ throughout, giving equal weight to the two signals.
The best-scoring synset is selected independently for the literal and
metaphorical meanings, so the same entry may (and often will) receive
different synsets for each.

%% -------------------------------------------------------------------------
\section{Running Example: \emph{Sapphire}}
\label{sec:example}

Table~\ref{tab:sapphire} illustrates the scoring for the entry
\emph{sapphire} under the theme \textsc{colour is mineral}.
The entry's literal meaning is \emph{transparent bright blue precious stone}
and its metaphorical meaning is \textsc{bright colour}; the source domain
is \{MINERAL\} and the target domain is \{COLOUR\}.

\begin{table}[ht]
\centering
\caption{Scores for all four \emph{sapphire} synsets under both methods
         ($\alpha=0.5$). Winning senses for each meaning are bold.}
\label{tab:sapphire}
\small
\begin{tabular}{llrrrr|rrrr}
\toprule
& & \multicolumn{4}{c|}{Literal (source: MINERAL)}
  & \multicolumn{4}{c}{Metaphorical (target: COLOUR)} \\
\cmidrule(lr){3-6}\cmidrule(lr){7-10}
Synset & Definition (abridged) & \multicolumn{2}{c}{Jac} & \multicolumn{2}{c|}{Emb}
                               & \multicolumn{2}{c}{Jac} & \multicolumn{2}{c}{Emb} \\
& & $\mathrm{sim}$ & total & $\mathrm{sim}$ & total
  & $\mathrm{sim}$ & total & $\mathrm{sim}$ & total \\
\midrule
\textbf{15019483-n} & precious transparent stone \ldots  & .444 & \textbf{.722} & .749 & \textbf{.875} & .000 & .000 & .363 & .182 \\
13372812-n & cut and polished sapphire \ldots           & .167 & .083           & .635 & .317           & .000 & .000 & .307 & .153 \\
\textbf{04969242-n} & a light shade of blue               & .143 & .071           & .575 & .288           & .000 & \textbf{.500} & .653 & \textbf{.827} \\
00383291-s & of colour of a blue sapphire (adj)          & .100 & .050           & .588 & .294           & .000 & .000 & .468 & .234 \\
\bottomrule
\end{tabular}
\end{table}

Several observations are worth noting.

\paragraph{Hypernym signal dominates for well-classified entries.}
For \texttt{omw-en-15019483-n} (\emph{precious transparent stone of rich
blue corundum}), the BFS reaches the synset \texttt{omw-en-14662574-n}
(\emph{mineral}), giving a hypernym score of~1.0.
For \texttt{omw-en-04969242-n} (\emph{a light shade of blue}), the BFS
reaches \texttt{omw-en-04959672-n} (\emph{chromatic colour}), again
scoring~1.0 against the target domain COLOUR.

\paragraph{Embeddings recover definition similarity that overlap misses.}
Under Jaccard, the metaphorical meaning \textsc{bright colour} scores 0.0
against \emph{a light shade of blue} because no content words overlap.
The embedding similarity is 0.653, correctly capturing that the two
phrases are semantically near-synonymous.
Without the hypernym signal this sense would still rank first for
metaphorical matching under embeddings (score 0.827 vs.\ 0.234 for the
next best), but the hypernym contribution of~0.5 alone is already
sufficient to select it correctly under Jaccard.

\paragraph{Correct senses are selected.}
Both methods agree on the correct pair:
\begin{itemize}
  \item \textbf{Literal}: \texttt{omw-en-15019483-n}
    (\emph{a precious transparent stone of rich blue corundum}),
    hypernym anchor: MINERAL $\to$ \texttt{omw-en-14662574-n}.
  \item \textbf{Metaphorical}: \texttt{omw-en-04969242-n}
    (\emph{a light shade of blue}),
    hypernym anchor: COLOUR $\to$ \texttt{omw-en-04959672-n}.
\end{itemize}

%% -------------------------------------------------------------------------
\section{Preliminary Results}
\label{sec:results}

We ran the full pipeline over all 8,747 thesaurus entries using the
embedding method ($\alpha = 0.5$, \texttt{embeddinggemma:latest},
batch size 64).
Table~\ref{tab:coverage} reports basic coverage statistics.

\begin{table}[ht]
\centering
\caption{Coverage of the full thesaurus run.}
\label{tab:coverage}
\begin{tabular}{lr}
\toprule
Metric & Count \\
\midrule
Total entries                  & 8,747 \\
Distinct headwords             & 6,880 \\
Entries with $\geq 1$ WN sense & 6,342 \\
Entries with 0 WN senses       & 2,405 \\
Unique synset texts embedded   & 31,610 \\
\midrule
Source domains (map)           &   213 \\
\quad of which: grouping nodes &    56 \\
\quad of which: single synsets &   157 \\
WN ancestry promotions         &    92 \\
\bottomrule
\end{tabular}
\end{table}

Of the 2,405 entries with no WordNet senses, most are idioms, fixed
phrases, or compound expressions whose headwords do not appear as lemmas
in OMW~2.0.
Threshold selection for accepting a match as reliable is an open
question: entries where both $\mathrm{sim}$ and $h$ are low should be
treated as unmatched.
We defer systematic evaluation against manually annotated gold data to
future work.

\subsection{Domain Coverage and Generalisation}
\label{subsec:domaincoverage}

We assess how well the domain map synsets cover the entries via two
matching strategies and compare them to the existing lemma-based approach:

\begin{description}
  \item[lemma] The current pipeline: the domain keyword appears as a
    lemma in some ancestor of the entry's literal synset.
  \item[exact] The domain's canonical synset (from the domain map) appears
    in the ancestor chain of the entry's literal synset.
  \item[+$k$ levels] The domain synset's ancestor at $k$ steps up the
    WordNet hierarchy appears in the entry's ancestor chain (i.e.\ we
    generalise the domain definition by $k$ WordNet levels).
\end{description}

\noindent
All percentages are computed over entries that have at least one WN sense.
Table~\ref{tab:domaincoverage} shows results for selected domains.

\begin{table}[ht]
\centering
\small
\caption{Coverage of selected source domains under different matching
strategies.  $N$ = entries in the domain's theme(s); no\_wn = fraction
without any WN sense; lemma, exact, +1, +2 = cumulative additional
entries matched (as \% of those with WN senses); miss = still unmatched
after +2 levels.}
\label{tab:domaincoverage}
\begin{tabular}{lrrrrrrr}
\toprule
Domain & $N$ & no\_wn & lemma & exact & +1 & +2 & miss \\
\midrule
\multicolumn{8}{l}{\textit{Well-matched concrete domains}} \\
\textsc{food}       & 123 & 25\% & 47\% & 45\% &  2\% &  1\% & 42\% \\
\textsc{body}       & 122 &  8\% & 50\% & 50\% &  0\% & 14\% & 24\% \\
\textsc{sound}      &  39 & 10\% & 54\% & 51\% &  5\% &  5\% & 31\% \\
\textsc{plant}      & 110 & 12\% & 48\% & 44\% &  6\% &  4\% & 39\% \\
\midrule
\multicolumn{8}{l}{\textit{Motivating example}} \\
\textsc{mineral}    &  46 & 13\% & 22\% & 20\% & 15\% &  2\% & 62\% \\
\midrule
\multicolumn{8}{l}{\textit{Domains gaining most from +1 generalisation}} \\
\textsc{amphibian}  &  17 & 11\% & 13\% & 13\% & 53\% &  0\% & 33\% \\
\textsc{waterbird}  &  19 & 15\% & 25\% & 25\% & 37\% &  0\% & 37\% \\
\textsc{implement}  &  37 & 21\% & 31\% & 27\% & 31\% &  0\% & 41\% \\
\textsc{contract}   &  17 & 17\% & 35\% & 28\% & 35\% & 0\% & 14\% \\
\textsc{building}   & 225 & 18\% &  9\% &  8\% & 22\% &  2\% & 62\% \\
\textsc{material}   &  67 & 32\% & 26\% & 26\% & 22\% &  6\% & 40\% \\
\midrule
\multicolumn{8}{l}{\textit{Domains gaining most from +2 generalisation}} \\
\textsc{fluid}      &  54 & 20\% & 13\% &  9\% & 11\% & 30\% & 41\% \\
\textsc{position}   &  92 & 23\% & 40\% & 30\% &  5\% & 25\% & 34\% \\
\textsc{vegetable}  &  22 & 18\% & 16\% &  5\% & 11\% & 27\% & 55\% \\
\midrule
\multicolumn{8}{l}{\textit{Abstract domains (hard to match via WordNet hierarchy)}} \\
\textsc{high}       & 259 & 26\% & 16\% &  7\% &  0\% &  0\% & 92\% \\
\textsc{distance}   & 195 & 34\% &  5\% &  3\% &  0\% &  0\% & 96\% \\
\textsc{heat}       & 120 & 35\% & 10\% &  7\% &  0\% &  2\% & 89\% \\
\midrule
\textbf{Total}      & 8,723 & 26\% & 24\% & 19\% &  4\% &  3\% & 67\% \\
\bottomrule
\end{tabular}
\end{table}

Several patterns emerge.
First, \textbf{lemma match and synset-exact are close overall} (24\% vs.\
19\%), validating the quality of the domain map's canonical synsets: the
synset-based approach is slightly more conservative but equally targeted.

Second, \textbf{one WordNet level of generalisation materially helps for a
subset of concrete domains}.
Notably, \textsc{amphibian} gains 53~percentage points at +1~level:
the canonical synset is too narrow (it covers only one order of
amphibians), but the parent synset covers the full class.
Similarly, \textsc{waterbird}, \textsc{implement}, and \textsc{building}
all gain more than 20~pp at +1~level.
This confirms the motivating observation that words like \emph{amber} and
\emph{gold} fail to match \textsc{mineral}'s canonical synset but succeed
at +1~level (covering the broader \emph{material} class).

Third, \textbf{abstract source domains remain largely uncovered at any
level}.
Domains such as \textsc{high}, \textsc{distance}, \textsc{heat}, and
\textsc{movement} have miss rates of 63--96\%, because most entries in
these themes denote abstract properties or manner adverbials that do not
sit in a WordNet hypernym chain passing through a recognisable domain
ancestor.
These represent a structural limit of the hypernym-matching approach and
point toward the need for distributional or frame-semantic signals for
such domains.

%% -------------------------------------------------------------------------
\section{Discussion}
\label{sec:discussion}

\paragraph{Strengths.}
The asymmetric use of hypernym chains — source domains for literal senses,
target domains for metaphorical senses — exploits the linguistic
structure of conceptual metaphor directly.
Where a clean hypernym path exists the signal is strong and interpretable,
and the matched ancestor synset provides a grounded semantic anchor.
The source domain taxonomy adds a second layer of structure: the 92
WordNet-inferred parent--child relationships among domains (e.g.\
\textsc{waterbird} $\to$ \textsc{bird} $\to$ \textsc{n\_top}) make the
ontology navigable and reusable for downstream tasks.

\paragraph{Domain map calibration.}
The coverage analysis (Section~\ref{subsec:domaincoverage}) reveals that
the domain map's canonical synsets are well-calibrated for concrete,
taxonomically organised domains (animals, body parts, food, plants) but
too narrow for several others.
Domains like \textsc{amphibian}, \textsc{waterbird}, and \textsc{implement}
would benefit from shifting their canonical synset one WordNet level up;
domains like \textsc{fluid} and \textsc{position} benefit more from
two levels.
The \textbf{lemma $\approx$ exact} pattern observed in the majority of
domains confirms that the lemma-based matching and the synset-based
approach are largely equivalent in practice, which validates using either
for downstream applications.

\paragraph{Limitations.}
\begin{itemize}
  \item \emph{Abstract source domains.}  Domains such as \textsc{high},
    \textsc{distance}, and \textsc{heat} have miss rates above 85\% at all
    levels tested.
    Their entries denote degree, scale, or property adjectives that do not
    follow a hypernym path through recognisable domain keywords.
    Alternative signals (frame-semantic, distributional, or corpus-based)
    are needed for this class of domain.
  \item \emph{Multi-synset domains.}  For 56 source domains (e.g.\
    \textsc{bird}) the canonical representation is split across several
    synsets, reflecting genuine polysemy or taxonomic ambiguity.
    The grouping node in the taxonomy has no synset of its own, so it
    cannot benefit from the WordNet ancestry promotion step.
  \item \emph{No POS filtering.}  We retrieve all synsets irrespective of
    part of speech because POS in the thesaurus reflects metaphorical usage.
    A post-hoc POS compatibility check could penalise implausible matches.
  \item \emph{Idioms and multi-word expressions.}  The 2,405 entries with
    no WordNet sense are overwhelmingly idioms, fixed phrases, or compound
    expressions absent from OMW~2.0.
  \item \emph{Threshold.}  Without gold-standard annotations the
    appropriate score threshold for accepting a match is unknown.
\end{itemize}

\paragraph{Future work.}
\begin{itemize}
  \item \emph{Adaptive domain synset generalisation.}
    Re-run hypernym matching using each domain's canonical synset (or its
    ancestor at the optimal WordNet level, as identified by the coverage
    analysis) instead of raw lemma string matching.
    This would directly address the \textsc{mineral}/\textsc{material}
    discrepancy observed for entries like \emph{amber} and \emph{gold}.
  \item \emph{Domain synset calibration.}
    For domains where +1 or +2 generalisation yields the largest gains
    (Table~\ref{tab:domaincoverage}), update the canonical synset in the
    domain map to the more general level and re-run to measure the effect
    on overall coverage.
  \item \emph{Abstract domain matching.}
    Investigate frame-semantic or distributional methods for domains
    (e.g.\ \textsc{high}, \textsc{distance}) that are not well served by
    WordNet hypernym chains.
  \item \emph{Manual evaluation.}
    Evaluate against a hand-annotated sample of 200--300 entries to
    establish precision and recall for sense selection.
  \item \emph{Cross-lingual projection.}
    Propagate synset links across languages via multilingual wordnets to
    extend \textsc{MetThes} coverage cross-lingually.
  \item \emph{Fine-tuned encoders.}
    Explore encoders trained on definition-similarity tasks to improve
    the definition similarity signal.
\end{itemize}

%% -------------------------------------------------------------------------
\section{Conclusion}
\label{sec:conclusion}

We have presented a method for automatically linking a large English
conceptual metaphor thesaurus to WordNet senses using two complementary
signals: definition similarity (Jaccard or embedding cosine) and hypernym
chain matching against the source and target domains encoded in each
theme's name.
The method correctly separates the literal and metaphorical senses of
entries such as \emph{sapphire} and produces interpretable, grounded
annotations for 6,342 of the 8,747 entries (72.5\%).
As a by-product we build a source domain taxonomy of 213 domains
organised as a shallow WordNet-anchored hierarchy, with 92 inter-domain
IS-A relationships inferred automatically from WordNet ancestry.

Coverage analysis shows that the existing lemma-based matching and the
new synset-based approach agree closely for most domains, but that a
targeted one-level generalisation of the domain synset could recover
substantial additional matches for concrete taxonomic domains
(e.g.\ +53~pp for \textsc{amphibian}, +37~pp for \textsc{waterbird}).
Abstract source domains remain a harder challenge, motivating future
work on frame-semantic and distributional methods.

The full pipeline is implemented in Python using the \texttt{wn} library
and the Ollama local inference server, with persistent embedding caching
to enable efficient reruns.

\bibliographystyle{plainnat}
\bibliography{paper}

\end{document}
